\subsection{Finite Massey vanishing and the retained integral lifting problem}

\begin{result}{P.14}{Known}
{The available Massey vanishing theorems do not supply the required compatible section}
Let $k$ be a number field and $p$ a prime.  Every defined $n$-fold
Massey product of classes in $H^1(k,\mathbb F_p(1))$, for $n\ge3$,
contains zero, with the successive products interpreted in their
appropriate Tate twists.  This is the rank-one coefficient case of
\citet[Theorem 7.4, author p.27]{RootHarpazWittenberg2023}.
For trivial coefficients and odd $p$, if $\zeta_p\notin k$, the
stronger condition that adjacent cups vanish suffices; it also
suffices over the maximal tame Galois quotient
\citep[Theorem 1.1]{RootMaireMinacRamakrishnaTan2026}.

These are finite-coefficient conclusions with an existential choice
of defining system.  They do not assert liftability of every specified
preceding section, simultaneous compatibility of all the obstruction
coordinates for the actual tripod action, or compatible integral
sections with one fixed ordinary pair.  In particular they do not
decide A.2, A.3, or the remainder in B.1.
\end{result}

\begin{proof}[Derivation and comparison]
We specify both the coefficient comparison and the quantifiers.
In Harpaz--Wittenberg's notation take
$N_i=\mathbb F_p(-i)$, $0\le i\le n$.  Then
\[
 M_{ij}=\operatorname{Hom}_{\mathbb F_p}(N_j,N_i)
       =\mathbb F_p(j-i),\qquad M_{i,i+1}=\mathbb F_p(1).
\]
Their Theorem 7.4 applies to these actual one-dimensional Galois
modules.  Proposition 7.3, author pp.26--27, identifies a defining
system with a continuous homomorphism to the corresponding finite
upper-triangular group modulo its top off-diagonal subgroup, retaining the
diagonal characters.  Vanishing means that \emph{some} defining
system lifts to the full triangular group.  It does not say that
the homomorphism belonging to an already specified defining system
lifts.  Taking every $N_i$ trivial gives their untwisted Theorem 1.3,
author p.3, as a special case.  The arithmetic fundamental-group
obstruction instead belongs to the particular predecessor in A.6;
its several coordinates share that predecessor.  The necessary
comparison to Massey expressions in P.1 retains these shared choices
and the Ihara action correction.  Separate existential choices of
defining systems need not be that common choice.

For the stronger 2026 theorem, the input consists of actual
characters $\chi_i:G_k\to\mathbb F_p$ with trivial coefficient
action, satisfying $\chi_i\cup\chi_{i+1}=0$.  Theorem 1.1 produces
a homomorphism into $U_{n+1}(\mathbb F_p)$ with those adjacent
entries.  Its Theorem 3.2 is a finite $p$-group embedding theorem
with local plans and additional tame ramification.  These source
and target hypotheses have a concrete mismatch with the proposed
application.  Over $\mathbb Q$, the residual cyclotomic character
has image $\mathbb F_p^\times$.  A finite arithmetic target retaining
that character has a quotient of order $p-1$, so for odd $p$ it is
not the $p$-group required by that theorem.  Restricting to
$\mathbb Q(\zeta_p)$ removes this particular obstruction but violates
the theorem's hypothesis $\zeta_p\notin k$.  At $p=2$ that hypothesis
already fails over every characteristic-zero field.  We do not
assert that a suitable twisted extension of the theorem is false;
it is an additional statement that would need proof.

There is also a relevant quantifier in Theorem 1.2 and its repeated
form, Theorem 4.2.  For a requested modulus $p^r$ the proof chooses
the additional ramified primes for the residual representation
with norms congruent to one modulo the exponent of
$U_{n+1}(\mathbb Z/p^r)$.  Thus the proof's choice of residual lift
can depend on $r$.  Part (i) gives a lift into the general linear
group; the unipotent lift in part (ii) additionally requires
$\zeta_{p^r}\in k_v$ at every prime ramified in the original
adjacent-entry representation.  The paragraph immediately after
the proof of Theorem 4.2 explicitly states that this construction
does not pass to an integral projective limit, because the required
auxiliary-prime congruence would have an infinite exponent.
No inverse system of integral tripod sections
with a retained ordinary abelian cocycle is supplied by that
statement.  Even liftability at each finite coefficient for a fixed
preceding datum would require the separate continuity and
compatibility analysis of A.20--A.24.

Finally, the auxiliary ramification cannot simply be inserted into
an integral correction of a fixed lower section.  Such a correction
in weight $w>1$ has coefficients $\mathbb Z_p(w)$.  At an auxiliary
prime $q\ne p$, its tame inertia value $v$ satisfies
\[
                     (q^{w-1}-1)v=0.
\]
The scalar is nonzero and the module is torsion-free, so $v=0$.
Modulo $p^a$ the same scalar can have a nonzero annihilator.  A
finite-coefficient construction using that annihilator is therefore
not an integral correction with the same lower data.  Allowing the
lower section and the pair to change introduces the old-support
matching equations proved in E.148 and E.159; none of the cited
theorems solves those equations along one infinite ordinary path.
This establishes the exact scope of the proposed reduction, without
turning its missing implication into a nonexistence theorem.
\end{proof}
